\documentclass{scrreprt}

% Dam0k1es

\setcounter{tocdepth}{3}                   % Setzt die Tiefe des Inhaltsverzeichnisses auf 3 Ebenen: Kapitel (chapter), Abschnitt (section), Unterabschnitt (subsection)
\usepackage{tabularx}                      % Ermöglicht erweiterte Tabellenfunktionen, insbesondere Tabellen mit automatisch anpassbarer Spaltenbreite

\usepackage[ngerman]{babel}                % Stellt Spracheinstellungen auf neues Deutsch um, beeinflusst z.B. Datumformat und Silbentrennung
\usepackage[german]{translator}            % Wird für die Übersetzung von Strings in Paketen verwendet, die nicht direkt von babel unterstützt werden

\usepackage[T1]{fontenc}                   % Wählt eine Schriftkodierung, die eine bessere Darstellung von Umlauten und anderen europäischen Zeichen ermöglicht
%\usepackage{microtype}                     % Verbessert den Textsatz durch Mikrotypografie

\usepackage{fixme}                         % Ermöglicht das Hinzufügen von Kommentaren und Notizen im Dokument, die leicht auffindbar sind
\fxsetup{status=draft}                     % Legt den Status des Dokuments fest; 'draft' bedeutet, dass Notizen sichtbar sind




% END Dam0k1es

\KOMAoptions{DIV      = 10,                % Gibt die Größe des Textbereichs an
             fontsize = 11pt,              % Schriftgröße
             paper    = a4,                % DIN A4-Papier
             parskip  = half}              % Kurzer Abstand statt eingerückter Absatzstart
\setlength{\headheight}{24pt}
\displaywidowpenalty=1000%                 % Leicht erhöhte Strafe für Witwen durch abgesetzten
                                           % Mathematiksatz
\widowpenalty=1000%                        % Leicht erhöhte Strafe für Witwen

% Auskommentiert aufgrund von 0x01
% \usepackage{iftex}                         % LuaLaTeX wird vorausgesetzt (für fontspec und selnolig)
% \RequireLuaTeX%

\usepackage{booktabs}                      % Bessere Tabellenlinien (eigene Befehle, Doku hilft)
\usepackage{csquotes}                      % Bessere Zitate (eigene Befehle, Doku hilft)
\usepackage{relsize}                       % Abgesetzte Zitate mit verkleinerter Schrift
\newenvironment{smallquotes}{\quote\smaller}{\endquote}
\SetBlockEnvironment{smallquotes}
\usepackage{ellipsis}                      % Korrigiert Platz nach \dots
\usepackage{etoolbox}                      % Um das ellipsis-Paket robust zu machen
\robustify\textellipsis%
\usepackage{enumerate}                     % Weitere Möglichkeiten für die enumerate-Umgebung
\usepackage{geometry}                      % Zentriert die Titelseite für die Druckausgabe
\usepackage{graphicx}                      % Einbinden von Bildern
\graphicspath{{Bilder/}}                      % Bilder können aus dem Verzeichnis img ohne Pfadangabe
                                           % eingebunden werden
\usepackage{icomma}                        % Korrigiert Platz um Kommas in Zahlen
\usepackage{listings}                      % Einbinden von Quelltexten
\usepackage{xcolor}                        % Nutzung von Farben
\lstset{%
    breaklines       = true,               % Automatischer Zeilenumbruch
    backgroundcolor  = \color{gray!8},
    commentstyle     = \textcolor{black!70},
    numbers          = left,
    numbersep        = 5pt,
    numberstyle      = \textcolor{black!40},
    showspaces       = false,
    showstringspaces = false,
    tabsize          = 4
}
\renewcommand*{\lstlistlistingname}{Listings-Verzeichnis}
\usepackage{mathtools}                     % Große Liste von Mathematikbefehlen, enthält amsmath
\usepackage{fontspec}                      % Für den normalen Schriftsatz (standardmäßig: lmodern)
\usepackage[final]{microtype}              % Verbessert die Lesbarkeit und verringert die Anzahl der
                                           % Trennungen
\usepackage{polyglossia}                   % Trennregeln für Deutsch und Englisch
\setmainlanguage[babelshorthands,          % Kürzel für Trennregeln wie im Paket babel
                 spelling = new]{german}
\setotherlanguage{english}
\usepackage[backend = biber,               % Standardmäßig ab TeXLive 2016
            sorting = none,                % Sortieren nach erster Nennung
            style   = ieee]{biblatex}      % Bibliographie im IEEE-Stil
\addbibresource{Quellen/Literatur.bib}
% Definiert den Bibliographiestil für Online-Quellen nach IEEE StyleGuide von 2009.
% Unter Umständen muss der Stil aktualisiert werden oder kann entfernt werden, wenn er in den
% BibLaTeX-Stil integriert wird.
% https://www.ieee.org/documents/ieeecitationref.pdf
\DefineBibliographyStrings{german}{%
    url = {\mkbibbrackets{Online}\adddot\addspace{}Verfügbar}
}
\renewbibmacro*{title}{%
    \ifboolexpr{
        test {\iffieldundef{title}}
    }
    {}
    {%
        \printtext[title]{%
            \printfield[sentencecase]{title}%
        }%
    }%
    \printfield{titleaddon}%
}
\DeclareBibliographyDriver{online}{%
    \usebibmacro{bibindex}%
    \usebibmacro{begentry}%
    \usebibmacro{author/editor+others/translator+others}%
    \adddot\addspace%
    \mkbibparens{\usebibmacro{date}}%
    \adddot\addspace%
    \usebibmacro{title}%
    \addspace%
    \usebibmacro{url}%
}
\usepackage{scrhack}                       % Verbessert die Kompatibilität zwischen KOMAskript und
                                           % einigen älteren Paketen
\usepackage[automark]{scrlayer-scrpage}    % Seitenlayout anpassen
\cohead*{}
\cofoot*{}
\rofoot*{\pagemark}
% Dieses Paket erzwingt LuaLatex. Damit muss der Font-Cache berechnet werden. 
% Dies führt zu einer viel langsameren Kompilierung des Projekts.   ID: 0x01
% \usepackage[ngerman]{selnolig}             % Unterdrückt selektiv Ligaturen
\usepackage[ngerman]{varioref}             % Referenziert Labels wo möglich und nötig mit
                                           % relativen Seitenangaben
\usepackage[breaklinks,                    % Trennung von Links erlauben
            colorlinks,                    % Links farbig setzen
            hyperfootnotes = false,        % Sehr fragil und unnötig
            linkcolor      = black,
            urlcolor       = blue]
           {hyperref}                      % Klickbare Hyperlinks
\usepackage[acronym,                       % Abkürzungsverzeichnis erstellen
            nogroupskip,                   % Kein Abstand zwischen Gruppen
            nonumberlist,                  % Keine nummerierte Liste der Einträge
            nopostdot,                     % Kein automatischer Punkt nach Einträgen
            xindy]
           {glossaries}                    % Zum Einbinden eines Abkürzungsverzeichnis und eines
                                           % Glossars
\GlsSetXdyCodePage{duden-utf8}             % Einstellungen für Xindy: Sortierung und UTF8
\setacronymstyle{short-long}               % Erste Abkürzung im Format kurz (lang)
\setglossarystyle{long}                    % Gleiche Breite für Spalten
\setlength\LTleft\parindent%               % Linke Spalte linksbündig statt zentriert
\setlength\LTright\fill%                   % Rechte Spalte füllt die restliche Seite
\makeglossaries%                           % Glossareinträge erstellen
\loadglsentries{Quellen/glossary}              % Glossareinträge aus dieser Datei lesen
\usepackage{cleveref}                      % Referenziert zusätzlich den Typ eines Labels

\AtBeginDocument{% Folgende Befehle am Beginn des Dokuments ausführen
    \hyphenation{Hash-al-go-rith-mus} % Hier globale Regeln zur Silbentrennung einfügen
    \renewcommand*{\acronymname}{Abkürzungsverzeichnis}
}

% Wenn URLs in der Bibliographie unbedingt an jeder Stelle umgebrochen werden können dürfen müssen,
% können die nachfolgenden Zeilen genutzt werden. Achtung: Sehr hässlich. Das \raggedright von unten
% kann in diesem Fall entfernt werden.
% \setcounter{biburlnumpenalty}{1}
% \setcounter{biburlucpenalty}{1}
% \setcounter{biburllcpenalty}{1}

\AtEndDocument{% Folgende Befehle am Ende des Dokuments ausführen
    \cleardoublepage%
    \raggedright%
    % \sloppypar% Kann als Alternative zum \raggedright eingesetzt werden. Blocksatz wie in Word.
    \printbibliography% Bibliographie ausgeben
}

\newcommand*{\Signature}{\Author\\[-0.3cm]\rule{6cm}{1pt}\\Autor}
\newcommand*{\SignatureZwei}{\AuthorZwei\\[-0.3cm]\rule{6cm}{1pt}\\Autor}

\newcommand*{\command}[1]{\texttt{#1}}  % Befehle innerhalb der Zeile in Schreibmaschinenschrift
                                        % setzen
\newcommand*{\definition}[1]{\emph{#1}} % Definitionen in kursiver Schrift setzen
\newcommand*{\filename}[1]{\texttt{#1}} % Dateinamen in Schreibmaschinenschrift setzen
\newcommand*{\programname}[1]{\gls{#1}} % Programme automatisch ins Glossar aufnehmen (nicht
                                        % unbedingt benötigt)
                                        
%\end{document}
 % Ausgelagerte Präambel

% Hier können die Informationen über die Ausarbeitung eingetragen werden
\newcommand*{\Course}{Projektarbeit} % Bei vorlesungsbegleitenden Ausarbeitungen
\newcommand*{\Subject}{Deepfakes und Social Engineering} % Thema der Ausarbeitung
\newcommand*{\Author}{Julian Faigle}
\newcommand*{\Matriculation}{86292}
\newcommand*{\Discipline}{ITS}
\newcommand*{\AuthorZwei}{Max Ernstschneider}
\newcommand*{\MatriculationZwei}{86464}
\newcommand*{\DisciplineZwei}{AIT}
\newcommand*{\Semester}{6}
\newcommand*{\Professor}{Prof. Roland Hellman}
\newcommand*{\Deadline}{15.08.2024}
\newcommand*{\Location}{Aalen}

% Um die Übersetzung zu beschleunigen, kann man hiermit nur Teile der Ausarbeitung übersetzen lassen
% \includeonly{FILENAME,FILENAME,...}

% Die folgenden Zeilen können für die Druckausgabe verwendet werden
% Achtung: Durch den vergrößerten Textbereich kann eine erneute Optimierung der Zeilenumbrüche
% notwendig sein.
% \KOMAoptions{BCOR     = 0mm,       % Hier die durchschnittliche Länge eintragen, die durch die
%                                    % Bindung verloren geht
%              DIV      = 11,        % Vergrößerter Textbereich für die Druckausgabe
%              headings = openright, % Kapitel starten immer auf der rechten Seite
%              twoside}
% \ofoot*{\pagemark}
% \ifoot*{}
% \renewcommand*{\Signature}{\rule{6cm}{1pt}\\Unterschrift}

% Metadaten für das PDF
\hypersetup{%
    pdfinfo={%
    Author={\Author\ (Matrikelnummer: \Matriculation); \AuthorZwei\ (Matrikelnummer: \MatriculationZwei)},
    Title={\Subject},
    Subject={\Course},
    Keywords={\Course, \Course, \Author, \Discipline, \Professor}
    }
}

\begin{document}
    \newgeometry{hmargin = {0cm, 0cm}} % Sicherstellen, dass die Titelseite zentriert ist
\begin{titlepage}
    \centering% Horizontal zentrieren
    \vspace*{\fill}% Vertikal zentrieren (zusammen mit unten)
    \linespread{1.2}
    \sffamily
    \textbf{\huge{\Course}}\\[1.2cm]
    \large{\textbf{\Subject}}\\[1.2cm]
    vorgelegt von\\[1.2cm]
    \Author\ (Matrikelnummer: \Matriculation)\\
    Studiengang \Discipline\\[1.0cm]
    \AuthorZwei\ (Matrikelnummer: \MatriculationZwei)\\
    Studiengang \DisciplineZwei\\[1.0cm]
    Semester \Semester\\[0.2cm]
    \includegraphics[scale=0.3]{Bilder/htw}\\[0.1cm]
    \textbf{Hochschule Aalen}\\[0.2cm]
    Hochschule für Technik und Wirtschaft\\[0.2cm]
    Betreut durch \Professor\\[0.2cm]
    \Deadline%
    \vspace*{\fill}% Vertikal zentrieren (zusammen mit oben)
\end{titlepage}
\restoregeometry% Bisheriges Seitenformat wiederherstellen
      % Titelseite
    \cleardoublepage%
\pagestyle{plain}
\pagenumbering{Roman}

\chapter*{Erklärung}
Wir versichern, dass wir die Ausarbeitung mit dem Thema \enquote{\Subject} selbstständig verfasst haben und keine anderen Quellen und Hilfsmittel als die angegebenen benutzt haben.
Die Stellen, die anderen Werken dem Wortlaut oder dem Sinn nach entnommen wurden, sind in jedem einzelnen Fall unter Angabe der Quelle als Entlehnung (Zitat) kenntlich gemacht worden.
Das Gleiche gilt für beigefügte Skizzen und Darstellungen.\\[4cm]
\parbox{\textwidth}{%
    \parbox[b]{7cm}{%
        \Location, den \today\\[-0.3cm]
        \rule{6cm}{1pt}\\
        Ort, Datum
    }
    \hfill%
    \parbox[b]{7cm}{%
        \Signature%
    }
}
\\[0.5cm]
\parbox{\textwidth}{%
    \hfill%
    \parbox[b]{7cm}{%
        \SignatureZwei%
    }
}
 % Eidesstattliche Erklärung

    % Hier eigene Kapitel einfügen

    \chapter{Kurzfassung}\label{ch:kurzfassung}

In den letzten Jahren haben technologische Fortschritte im Bereich der \gls{ki} und des maschinellen Lernens zu einer signifikanten Entwicklung in der digitalen Manipulation von Medien geführt.
Insbesondere Deepfakes, bei denen KI verwendet wird, um realistisch aussehende Bilder, Videos oder Audiodateien zu erstellen, die es so nie gegeben hat, haben die Aufmerksamkeit von Forschern, Sicherheitsexperten und der breiten Öffentlichkeit gleichermaßen erregt.
Diese Technologie, die ursprünglich zur Verbesserung von visuellen Effekten und für kreative Zwecke entwickelt wurde, hat sich mittlerweile zu einem Werkzeug entwickelt, welches auch für bösartige Zwecke missbraucht werden kann.

Deepfakes basieren auf komplexen Algorithmen wie Generative Adversarial Networks (GANs), die es ermöglichen, Gesichter in Videos nahtlos auszutauschen, Stimmen zu imitieren oder Personen in Szenarien darzustellen, die nie stattgefunden haben.
Diese Fähigkeit, die Realität auf eine Weise zu verzerren, die für das menschliche Auge oft schwer zu erkennen ist, bringt erhebliche Risiken für die Gesellschaft.
Von der Verbreitung von Fehlinformationen über politische Manipulationen bis hin zu Erpressungen und Identitätsdiebstahl – die Einsatzmöglichkeiten von Deepfakes sind vielfältig und gefährlich.

Gleichzeitig hat Social Engineering, die Kunst der Manipulation von Menschen, um vertrauliche Informationen zu erhalten oder bestimmte Aktionen auszulösen, eine neue Dimension erreicht.
Social Engineering nutzt gezielt psychologische Techniken, um das Vertrauen von Individuen zu gewinnen und sie dazu zu bringen, sicherheitsrelevante Fehler zu begehen.
In einer zunehmend digitalisierten Welt, in der persönliche und berufliche Interaktionen häufig online stattfinden, wird die Verbindung zwischen Social Engineering und digitalen Technologien immer enger.

Die Kombination von Deepfakes und Social Engineering stellt eine besonders gefährliche Bedrohung dar.
Angreifer können Deepfakes nutzen, um das Vertrauen ihrer Zielpersonen zu erschleichen oder deren Entscheidungen zu beeinflussen, indem sie manipulierte Inhalte einsetzen, die beispielsweise vertraute Gesichter oder Stimmen imitieren.
Solche Angriffe sind schwer zu erkennen und zu bekämpfen, da sie sowohl technisches Wissen als auch ein tiefes Verständnis menschlicher Verhaltensmuster erfordern.

In dieser Ausarbeitung werden verschiedene Technologien im Bereich Video- und Audio-Deepfakes beleuchtet und vorgestellt.
Sie werden in den Kontext von Social Engineering eingeordnet und praktisch angewendet, um beispielhafte Szenarien darzustellen. % Kurzfassung

    {% Änderungen lokal für Verzeichnisse halten
    \renewcommand*{\chapterpagestyle}{plain}
    \microtypesetup{protrusion = false} % Unschön für nummerierte Verzeichnisse
    \tableofcontents%
    % Weitere Verzeichnisse können hier bei Bedarf hinzugefügt werden
    % \listoffigures%
    % \listoftables%
    % \lstlistoflistings%
    \microtypesetup{protrusion = true}
    \printglossary[type=\acronymtype] % Abkürzungen separat ausgeben
    \printglossary% Glossar ausgeben
    \cleardoublepage%
}

% Ab jetzt anderer Stil für Kapitel
\renewcommand*{\chapterformat}{\LARGE\sffamily\bfseries\textup\thechapter.\enskip}
\setkomafont{chapter}{\LARGE\sffamily\bfseries\textup}
\renewcommand*{\chapterheadstartvskip}{\vspace{-4.5em}}
\renewcommand*{\chapterheadendvskip}{\vspace{-1.4em}\rule{\textwidth}{1pt}\vspace{20pt}}

\pagestyle{plain}
\pagenumbering{arabic}
    % Verzeichnisse, Vorbereitung Layout eigene Kapitel

    % \input{}
\chapter{Grundlagen}\label{ch:grundlagen}
\input{Inhalt/Grundlagen-Einführung_in_Deepfakes}
\section{Social Engineering}\label{sec:social_engineering}
Unter Social Engineering werden alle Angriffe zusammengefasst, die die Schwachstelle Mensch ausnutzen.
Es wird durch verschiedene Techniken versucht, an private oder sensible Inhalte von Personen zu gelangen.
Social Engineering ist heutzutage eine der größten Gefahren im digitalen Raum.
Kryptografische Verfahren und Protokolle wurden über die Jahre immer besser.
Ist ein System bzw. eine digitale Infrastruktur richtig gehärtet, sind Angriffe wie Brute-Force oder Dictionary-Attacks wirkungslos.
Außerdem werden durch moderne \glspl{ids} sowie \gls{threat-intelligence} technische Angriffe immer schneller erkannt und blockiert.
Laut einem Paper aus 2018 sind 84\% aller Cyber-Angriffe auf Social Engineering zurückzuführen.
Zudem haben Social Engineering Angriffe eine höhere Erfolgschance als herkömmliche Methoden.
Laut des \gls{dbir} von Verizon waren 2024 $45\%$ der erfassten Cyberangriffe auf Social Engineering zurückzuführen.
Bei $~83\%$ der 3661 reporteten Incidents wurden Daten extrahiert.\cite{verizon-2024-dbir}

% Kommentar: Es wäre hilfreich, ein konkretes Beispiel für einen erfolgreichen Social Engineering Angriff einzufügen, um die Einführung lebendiger und verständlicher zu machen.

\subsection{Verschiedene Typen von Social Engineering}\label{subsec:verschiedene-typen-von-social-engineering}
Social Engineering lässt sich in verschiedene Bereiche untergliedern, die folgenden Beschreibungen richten sich nach dem Artikel: ``Social Engineering Attacks: A Survey''\cite{social-engineering-a-survey}

Hier werden Social Engineering Angriffe in zwei Kategorien unterteilt.
\begin{itemize}
    \item Human-based
\end{itemize}
    Diese Angriffe werden manuell von einem Menschen ausgeführt.
    Sie sind in der Regel spezifisch auf das Opfer angepasst und mit höherem Aufwand verbunden.
    Dafür sind die Erfolgschancen höher als bei automatisierten Angriffen.
\begin{itemize}
    \item Computer-based
\end{itemize}
Diese Angriffe werden automatisiert durchgeführt.
Sie sind qualitativ deutlich schlechter als ihr Gegenstück, dafür werden sie in hoher Quantität durchgeführt.
Hierzu zählen Phishing-Mails oder SMS.
Es gibt verschiedene Tools, um solche Angriffe durchzuführen, ein bekanntes ist das ``\href{https://github.com/trustedsec/social-engineer-toolkit}{Social Engineering Toolkit}''.

Des Weiteren können Social Engineering Angriffe in drei weitere Kategorien unterteilt werden.
\begin{itemize}
    \item Social-based
\end{itemize}
Diese Form von Social Engineering Angriffen besteht aus zwischenmenschlicher Interaktion.
Dabei spielt sie mit der Psychologie und den Emotionen der Zielperson.
Diese Form von Social Engineering birgt ein hohes Risiko, hat aber ebenfalls eine hohe Erfolgschance, da der Angreifer im direkten oder indirekten Kontakt mit dem Opfer steht.
Beispiele hierfür wären: Baiting, Spear-Phishing, aber auch Dinge wie der \gls{enkeltrick}.

\begin{itemize}
    \item Technical-based
\end{itemize}
Hier werden Angriffe übers Internet remote ausgeführt.
Dafür werden Social-Media Plattformen und Online-Dienste verwendet, um Passwörter, Kreditkarteninformationen oder personenbezogene Daten zu stehlen.
Hierzu zählen zum Beispiel Phishing-Kampagnen oder gefälschte Webseiten.

\begin{itemize}
    \item Physical-based
\end{itemize}
Physical-based Angriffe geschehen abseits des Internets in der realen Welt.
Dabei werden durch physisches Handeln Informationen erschlossen.
Ein Beispiel wäre das Durchsuchen von Müllcontainern (auch Dumpster-Diving genannt) nach sensiblen Dokumenten.

% Kommentar: Es wäre nützlich, Beispiele für Gegenmaßnahmen gegen die einzelnen Typen von Angriffen zu nennen, um die Informationen greifbarer zu machen.

Je nachdem, aus welchem Blickwinkel die verschiedenen Techniken des Social Engineerings betrachtet werden, können diese in noch mehr verschiedene Kategorien eingeteilt werden.
Neben Human-, Computer-, Social-, Technical- und Physical-based Social Engineering ist die zusätzliche Unterscheidung in \textbf{direkt} und \textbf{indirekt} sinnvoll.
Ersteres benötigt direkten Kontakt zwischen Angreifer und Opfer, dabei zählen physischer Kontakt sowie Telefonate.
Beispiele sind: physical access, shoulder surfing, dumpster diving, phone social engineering, pretexting, impersonation on help desk call.
Indirekte Angriffe sind entsprechend analog dazu.
Hierzu zählen: phishing, fake software, Pop-Up windows, ransomware, SMSishing, online social engineering.

\subsection{Überblick über gängige Angriffe}\label{subsec:uberblick-uber-gangige-angriffe}
Phishing ist eine der am weitesten verbreiteten Social Engineering-Techniken.
Ziel dieser Angriffe ist es, private oder vertrauliche Daten der Opfer zu stehlen.
Dabei werden hauptsächlich E-Mails, SMS, Anrufe oder gefälschte Webseiten eingesetzt, um die Opfer zur Preisgabe ihrer Informationen zu verleiten.
Phishing lässt sich grob in folgende Kategorien unterteilen\cite{advanced-social-engineering-attacks}:
\begin{itemize}
    \item \textbf{Spear Phishing}: Diese Methode ist zielspezifisch und verwendet oft durch Open Source Intelligence (OSINT) gesammelte Informationen, um maßgeschneiderte E-Mails zu erstellen.
     Die Nachrichten wirken dadurch besonders glaubwürdig und erhöhen die Erfolgschancen des Angriffs.
    \item \textbf{Whaling Phishing}: Hierbei handelt es sich um Angriffe auf hochrangige Ziele, wie Führungskräfte oder Personen in Schlüsselpositionen.
     Diese Angriffe sind oft sehr aufwendig und spezifisch auf das Ziel zugeschnitten, um wertvolle Informationen zu erlangen.
    \item \textbf{Vishing}: Voice Phishing, bei dem Telefonanrufe oder Sprachdienste wie Teams genutzt werden, um sensible Informationen zu erlangen.
     Die Angreifer geben sich häufig als vertrauenswürdige Institutionen oder Personen aus, um das Vertrauen des Opfers zu gewinnen.
\end{itemize}

Eine weitere Technik ist \textbf{Baiting}.
Baiting, auch als Road Apples bekannt, verleitet Personen dazu, auf etwas zu klicken oder ein Gerät zu benutzen, um vermeintlich etwas gratis zu erhalten.
Ein bekanntes Beispiel hierfür sind E-Mails mit einem Gewinn, für den man sich nur noch registrieren braucht, um ihn zu erhalten.
Außerdem gehören auch infizierte USB-Sticks, die in der Hoffnung verteilt werden, dass jemand sie benutzt, zu Baiting dazu.
Bei Bad-USBs wird auf die Neugierde des Menschen gesetzt.
Durch das Einstecken des USB-Sticks in einen Computer kann Schadsoftware installiert werden, die es den Angreifern ermöglicht, auf das System zuzugreifen.

\textbf{Tailgating Attacks} beziehen sich auf das unerlaubte Verschaffen von Zutritt zu gesicherten Bereichen, indem zum Beispiel einer autorisierten Person gefolgt wird.
Auch Angriffe auf die Sicherheitsmechanismen, wie z.B. das Kopieren eines NFC- oder RFID-Tags gehören in diese Kategorie.
Solche Angriffe ermöglichen es dem Angreifer, physisch gesicherte Bereiche zu betreten und dort Informationen zu stehlen oder Schaden anzurichten.

\subsection{Gegenmaßnahmen gegen Social Engineering Angriffe}\label{subsec:gegenmassnahmen-gegen-social-engineering-angriffe}
Die Abwehr von Social Engineering Angriffen erfordert eine Kombination aus präventiven und reaktiven Maßnahmen.
Eine der effektivsten Präventionsstrategien ist die \textbf{Schulung der Mitarbeiter}.
Durch regelmäßige Schulungen und Sensibilisierungsprogramme können Mitarbeiter lernen, die Anzeichen von Social Engineering Angriffen zu erkennen und angemessen darauf zu reagieren.
Dies umfasst das Überprüfen der Authentizität und Integrität von Nachrichten, sei es per E-Mail, SMS oder Telefon.

\textbf{Überprüfung der Authentizität und Integrität von Nachrichten}: Mitarbeiter sollten stets darauf achten, ungewöhnliche oder verdächtige Nachrichten sorgfältig zu prüfen.
Dazu gehört, den Absender zu überprüfen, auf Rechtschreib- und Grammatikfehler zu achten und Links sowie Anhänge nicht ohne weiteres zu öffnen.
Wenn Zweifel bestehen, sollte die Nachricht direkt beim vermeintlichen Absender verifiziert werden.

Da human-basierte Social Engineering Angriffe schwer oder gar nicht automatisiert zu erkennen sind, ist die \textbf{Schadensbegrenzung} von entscheidender Bedeutung.
Hier kommen verschiedene technische und organisatorische Maßnahmen ins Spiel:

\begin{itemize}
    \item \textbf{Domain-Tiering}: Diese Technik hilft, die Auswirkungen eines erfolgreichen Angriffs zu minimieren, indem unterschiedliche Sicherheitsstufen für verschiedene Domänen innerhalb eines Unternehmens festgelegt werden.
     Dadurch kann ein kompromittierter Bereich isoliert und der Schaden begrenzt werden.
    \item \textbf{Notfallmanagement}: Ein effektives Notfallmanagement umfasst klare Protokolle und Verantwortlichkeiten für den Fall eines Angriffs.
     Regelmäßige Schulungen und Übungen stellen sicher, dass alle Mitarbeiter wissen, wie sie im Ernstfall reagieren müssen.
     Dies beinhaltet auch die schnelle Identifikation und Isolation kompromittierter Systeme sowie die Benachrichtigung betroffener Personen und Behörden.
\end{itemize}

Zusätzlich zu diesen Maßnahmen können technische Hilfsmittel den Schutz vor Social Engineering Angriffen verbessern:

\begin{itemize}
    \item \textbf{E-Mail-Sicherheitslösungen}: Tools wie E-Mail-Filter und Anti-Phishing-Software können verdächtige Nachrichten erkennen und blockieren, bevor sie die Mitarbeiter erreichen.
    \item \textbf{Zwei-Faktor-Authentifizierung (2FA)}: Durch die Implementierung von 2FA wird ein zusätzlicher Schutzlayer hinzugefügt, der es Angreifern erschwert, Zugang zu sensiblen Systemen und Daten zu erlangen, selbst wenn sie die Anmeldedaten eines Mitarbeiters gestohlen haben.
    \item \textbf{Netzwerküberwachung und \glspl{ids}}: Diese Systeme überwachen den Netzwerkverkehr auf verdächtige Aktivitäten und können Angriffe frühzeitig erkennen und abwehren.
\end{itemize}

Ein ganzheitlicher Ansatz, der sowohl präventive als auch reaktive Maßnahmen umfasst, ist unerlässlich, um die Widerstandsfähigkeit eines Unternehmens gegenüber Social Engineering Angriffen zu erhöhen.
Durch die Kombination aus regelmäßiger Mitarbeiterschulung, technischer Absicherung und einem robusten Notfallmanagement kann das Risiko solcher Angriffe erheblich reduziert werden.\cite{social-engineering-a-survey, bsi-social-engineering}

% Kommentar: Es wäre sinnvoll, abschließend auf die Bedeutung der kontinuierlichen Überprüfung und Anpassung der Maßnahmen hinzuweisen, um auf neue Bedrohungen und Techniken der Angreifer reagieren zu können.

\section{Bekannte Social Engineering Angriffe}\label{sec:bekannte-social-engineering-angriffe}
\subsection*{Reenacted Video Call in Hong Kong}
Ein besonders eindrucksvolles Beispiel für die Nutzung von Social Engineering in Verbindung mit Deepfake-Technologie
ereignete sich Anfang des Jahres bei einem Unternehmen in Hongkong.
Ein Finanzmitarbeiter wurde von Betrügern dazu gebracht, \$25 Millionen zu überweisen.
Wie die Hongkonger Polizei berichtet, wurde der Mitarbeiter zu einem Videoanruf eingeladen, jede Person in diesem
Meeting war jedoch eine Deepfake. Mittels Echtzeit Face- und Voice-Reenactment wurde das gesamte Meeting gefälscht.

Der Betrug begann mit einer Nachricht, die angeblich vom Finanzchef des Unternehmens in Großbritannien stammte und von einer geheimen Transaktion sprach. Obwohl der Mitarbeiter zunächst misstrauisch war und einen Phishing-Versuch vermutete, ließ er sich durch die anschließende Videokonferenz überzeugen, da die anwesenden Personen wie seine tatsächlichen Kollegen aussahen und klangen.

Dieser Fall zeigt eindrucksvoll, wie fortschrittlich und gefährlich Deepfake-Technologie inzwischen ist. Sie wurde
hier nicht nur verwendet, um ein realistisches Video der Zielperson zu erstellen, sondern auch, um mehrere scheinbar
authentische Teilnehmer an einem Videoanruf zu simulieren. Die Täuschung flog erst auf, als der Mitarbeiter
nachträglich beim Hauptsitz des Unternehmens nachfragte\cite{cnn-deepfake-scam-hong-kong}.

\subsection*{Marriott-Hotel Databreach}
Ein weiteres Beispiel für einen Social Engineering Angriff ereignete sich im Juni 2022 im Marriott-Hotel am Flughafen von Baltimore im US-Bundesstaat Maryland. Kriminelle hatten sich mittels Social Engineering durch einen Mitarbeiter des Hotels Zugang zum Netzwerk verschafft und 20 GB an Daten abgeschöpft, darunter auch Kreditkartendaten von Gästen und interne Geschäftsdaten des Hotels. Marriott hat die Strafverfolgungsbehörden eingeschaltet und wird nach eigenen Abgaben etwa 400 Personen benachrichtigen. Eine Lösegeldforderung der Erpresser lehnte die Hotelkette ab\cite{marriott-data-breach}.

\subsection*{Phishing-Kampanie United States Department of Labor}
Ein weiteres Beispiel für einen Phishing-Angriff ist eine Kampagne aus dem Jahre 2022, die die United States Department
of Labor (DoL) imitiert und Empfänger auffordert, Angebote einzureichen, um Office 365-Anmeldeinformationen zu stehlen.
Die E-Mails passierten gekaperte Server, die gemeinnützigen Organisationen gehören, um E-Mail-Sicherheitsblöcke zu
umgehen. Außerdem wurde die Sender Adresse gespooft um den tatsächlichen Domains des DoL zu entsprechen.

Die Angreifer geben sich als leitender DoL-Mitarbeiter aus, der den Empfänger einlädt, sein Angebot für ein laufendes
Regierungsprojekt einzureichen. Die E-Mails enthalten einen gültigen Briefkopf, professionell gestalteten Inhalt und
einen dreiseitigen PDF-Anhang mit einem scheinbar legitimen Formular.

Das PDF enthält eine „BID“-Schaltfläche, die die Opfer auf eine Phishing-Seite weiterleitet.
Die gefälschte Seite sieht überzeugend aus und enthält identisches HTML und CSS wie die echte. Die Bedrohungsakteure
haben sogar eine Anweisungs-Pop-up-Nachricht hinzugefügt, um die Opfer durch (Phishing-)Prozess zu führen.

Wer für ein Projekt bietet, wird zu einem Formular zur Erfassung von Anmeldeinformationen weitergeleitet, das die
E-Mail-Adresse und das Passwort von Microsoft Office 365 der Opfer abfragt\cite{phishing-dol}.


\chapter{Deepfake Varianten}\label{ch:deepfake-varianten}
\section{Video-Deepfakes}\label{sec:video-deepfakes2}

Video-Deepfakes sind eine besondere Form von gefälschten Medieninhalten, bei denen Gesichtsbilder in Videos so manipuliert werden, dass sie authentisch erscheinen. Diese Technologie hat erhebliche Auswirkungen auf die Bereiche Sicherheit, Verifikation und Social Engineering. Im Folgenden werden verschiedene Methoden zur Erstellung von Video-Deepfakes beschrieben und ihre technischen Details, Unterschiede, Gemeinsamkeiten sowie Vor- und Nachteile diskutiert.

\subsection{Veraltete Methode: 3D Modellierung}\label{subsec:3d-modellierung}

Die frühe Methode zur Erstellung von Video-Deepfakes basierte auf der 3D-Modellierung. Hierbei wird ein dreidimensionales Modell des Gesichts einer Person erstellt, das anschließend animiert und in ein Video eingefügt wird.

\textbf{Technische Details:}
\begin{itemize}
    \item Erstellung eines 3D-Gesichtsmodells durch Scannen oder Fotogrammetrie.
    \item Animierung des Modells basierend auf Zielausdrücken oder Bewegungsdaten.
    \item Einfügung des animierten Modells in das Zielvideo mithilfe von Tracking- und Rendering-Techniken.
\end{itemize}

\textbf{Vor- und Nachteile:}
\begin{itemize}
    \item \textbf{Vorteile:} Hohe Kontrolle über die Gesichtsausdrücke und -bewegungen, gute Qualität bei statischen Szenen.
    \item \textbf{Nachteile:} Hoher Aufwand bei der Modellierung und Animation, Schwierigkeiten bei der realistischen Darstellung von dynamischen Szenen und feinen Details \cite{Deepfakes-An-Overview}.
\end{itemize}

\textbf{Einsatz im Kontext von Social Engineering:} Aufgrund des hohen Aufwands und der technischen Expertise, die für die Erstellung benötigt wird, sind 3D-Modellierungs-Deepfakes weniger verbreitet in Social Engineering-Angriffen.
Außerdem müssen die neuronalen Netze auf jedes Ziel- sowie Quellgesicht mit einigen Daten (~Minuten Video) neu trainiert werden.
Deshalb und aus Gründen der Performance ist dieser Ansatz für Echtzeitanwendungen, wie Videoanrufen, nicht geeignet \cite{social-engineering-a-survey}.

\subsection{Generative Adversarial Networks (\glspl{gan})}\label{subsec:gans}

Mit dem Aufkommen von Generative Adversarial Networks (\glspl{gan}) hat sich die Erstellung von Deepfakes erheblich vereinfacht und verbessert.
\glspl{gan} bestehen aus zwei neuronalen Netzwerken, einem Generator und einem Diskriminator, die gegeneinander trainiert werden.\cite{Deepfakes-a-survey-and-introduction-to-the-topical-collection}

\textbf{Technische Details:}
\begin{itemize}
    \item Der \textbf{Generator} erzeugt Bilder, die versuchen, echte Bilder nachzuahmen.
    \item Der \textbf{Diskriminator} bewertet die Bilder des Generators und unterscheidet zwischen echten und generierten Bildern.
    \item Durch diesen Wettbewerb verbessert sich die Qualität der erzeugten Bilder stetig.
\end{itemize}

\textbf{Vor- und Nachteile:}
\begin{itemize}
    \item \textbf{Vorteile:} Hohe Qualität der erzeugten Bilder, Möglichkeit zur Erstellung realistischer und dynamischer Gesichtsausdrücke.
    \item \textbf{Nachteile:} Erfordert große Datenmengen und Rechenressourcen für das Training, anfällig für Artefakte und Unstimmigkeiten bei komplexen Bewegungen \cite{Deepfakes-An-Overview}.
\end{itemize}

\textbf{Einsatz im Kontext von Social Engineering:} \gls{gan}-basierte Deepfakes sind effektiver und leichter zu erstellen, was sie zu einem mächtigen Werkzeug für Angreifer im Bereich Social Engineering macht.
Sie können verwendet werden, um gefälschte Videos zu erstellen, die Vertrauen erwecken und die Opfer täuschen.
\glspl{gan} sind ebenfalls zu ineffizient in der Implementierung, um in Echtzeitanwendungen Gebrauch zu finden.

\subsection{\gls{fsgan} und FSGANv2}\label{subsec:fsgan}

Eine Weiterentwicklung der \gls{gan}-Technologie sind die Face Swapping \glspl{gan} (\gls{fsgan}) und ihre verbesserte Version, FSGANv2.
Diese Technologien sind speziell für das Gesichtstausch und die Gesichtsnachstellung entwickelt worden.

\textbf{Technische Details:}
\begin{itemize}
    \item \gls{fsgan} nutzt \glspl{gan}, um Gesichtszüge von einer Quelle auf ein Zielvideo zu übertragen, ohne dass eine explizite 3D-Modellierung erforderlich ist.
    \item FSGANv2 verbessert diese Methode durch bessere Algorithmen zur Anpassung der Gesichtsausdrücke und -bewegungen sowie durch die Verwendung von fortschrittlichen Netzwerktechniken, um die Konsistenz und Realismus zu erhöhen \cite{fsganv2}.
\end{itemize}

\textbf{Vor- und Nachteile:}
\begin{itemize}
    \item \textbf{Vorteile:} Sehr realistische Ergebnisse, weniger Training und Daten erforderlich im Vergleich zu reinen \glspl{gan}, bessere Anpassung an unterschiedliche Gesichtsausdrücke und Beleuchtungen.
    \item \textbf{Nachteile:} Trotz Verbesserungen immer noch anfällig für subtile Unstimmigkeiten, die bei genauer Betrachtung auffallen können \cite{face-swapping-and-reenactment}.
\end{itemize}

\textbf{Einsatz im Kontext von Social Engineering:} \gls{fsgan} und FSGANv2 sind äußerst effektiv für Social Engineering-Angriffe, da sie hochrealistische Deepfakes erstellen können, die schwer zu erkennen sind.
Sie können verwendet werden, um falsche Identitäten zu erstellen und Vertrauen zu gewinnen, was das Risiko und den potenziellen Schaden solcher Angriffe erhöht \cite{deepfacelab}.

\subsection{Gemeinsamkeiten und Unterschiede}\label{subsec:gemeinsamkeiten-unterschiede}

\textbf{Gemeinsamkeiten:}
\begin{itemize}
    \item Alle Methoden zielen darauf ab, realistische Fälschungen zu erstellen, die schwer zu erkennen sind.
    \item Nutzung von KI und maschinellem Lernen zur Verbesserung der Qualität und Realismus der erzeugten Videos.
\end{itemize}

\textbf{Unterschiede:}
\begin{itemize}
    \item 3D-Modellierung erfordert manuelle Arbeit und ist weniger flexibel bei dynamischen Szenen.
    \item \glspl{gan} und deren Weiterentwicklungen (\gls{fsgan}, FSGANv2) bieten automatisierte Lösungen mit höherer Qualität und Flexibilität.
    \item \gls{fsgan} und FSGANv2 sind speziell auf Gesichtsmanipulation optimiert und bieten bessere Ergebnisse bei der Anpassung an verschiedene Bedingungen \cite{Deepfakes-a-survey-and-introduction-to-the-topical-collection}.
\end{itemize}

Insgesamt zeigt sich, dass die Weiterentwicklung der Deepfake-Technologien, insbesondere durch \glspl{gan} und deren Spezialisierungen wie \gls{fsgan} und FSGANv2, erheblich zur Verbesserung der Qualität und Realismus beigetragen hat.
Dies stellt jedoch auch eine größere Bedrohung im Bereich des Social Engineering dar, da die Täuschungsabsicht hinter den erzeugten Medieninhalten immer schwerer zu durchschauen ist.


\section{Face Swapping vs. Reenactment}\label{sec:face-swapping-vs.-reenactment}
Es gibt verschiedene Techniken von Video Deepfakes, im Folgenden werden Face-Swapping, sowie Reenactment näher
betrachtet.

\section{Face Swapping}\label{sec:face-swapping}
Face Swapping, eine weit verbreitete Technik innerhalb der Deepfake-Technologie, beinhaltet das Austauschen eines
Gesichts in einem Bild oder Video durch das Gesicht einer anderen Person.
Diese Methode hat insbesondere in den letzten Jahren erhebliche Fortschritte gemacht, vor allem durch die Entwicklung von \glspl{gan}.
Bei Face Swapping wird das Gesicht der Zielperson durch ein anderes Gesicht ersetzt, wobei Merkmale wie Hautfarbe,
Beleuchtung und Gesichtsausdrücke so angepasst werden, dass das Ergebnis möglichst realistisch wirkt.

Im Rahmen des \glspl{fsgan}\cite{face-swapping-and-reenactment} Ansatzes wird ein neuartiges rekurrentes neuronales Netzwerk (RNN) für das Face Swapping vorgestellt, das sowohl Pose- als auch Ausdrucksvariationen anpasst und auf ein einzelnes Bild oder eine Video-Sequenz angewendet werden kann.
Die kontinuierliche Interpolation der Gesichtsansichten erfolgt auf der Grundlage von Delaunay-Triangulation und baryzentrischen Koordinaten.
Ein Gesichtsvollständigungsnetzwerk behandelt verdeckte Gesichtsregionen, während ein Gesichtsmischungsnetzwerk für das nahtlose Einfügen der beiden Gesichter sorgt und dabei die Hautfarbe und Lichtverhältnisse des Zielbildes bewahrt.
Diese Methode übertrifft sowohl qualitativ als auch quantitativ bestehende Systeme.

Frühere Methoden für Face Swapping basierten oft auf 3D-Gesichtsmodellen, die entweder geschätzt oder festgelegt wurden.
Diese Modelle wurden dann an die Eingabebilder angepasst und als Proxy für das Übertragen von Intensitäten (Swapping) oder das Steuern von Gesichtsausdrücken und -ansichten (Reenactment) verwendet.
Neuere Ansätze nutzen tiefe Netzwerke, insbesondere \glspl{gan}, um realistische Bilder von gefälschten Gesichtern zu erzeugen.
Diese Methoden sind jedoch oft auf spezifische Subjekte trainiert und erfordern umfangreiche Datensätze, um akzeptable Ergebnisse zu erzielen.
Im Gegensatz dazu ist \gls{fsgan} subjektsagnostisch und benötigt kein training auf spezifischen Gesichtern, was seine Anwendung deutlich vereinfacht und erweitert.

\section{Reenactment}\label{sec:reenactment}
Reenactment, auch als Face Transfer oder Puppeteering bekannt, ist eine Technik, bei der die Gesichtsausdrücke und -bewegungen eines Ausgangsbildes oder -videos auf ein Zielbild oder -video übertragen werden.
Dies ermöglicht es, das Gesicht der Zielperson so zu manipulieren, dass es die gleichen Bewegungen und Ausdrücke wie das Ausgangsgesicht zeigt.
Diese Technik wird häufig in der Unterhaltungsindustrie, aber auch in Bereichen wie Datenschutz und der Erstellung von Trainingsdaten verwendet.

Die \gls{fsgan}-Methode bietet eine innovative Lösung für das Reenactment, die keine spezifischen Trainingsdaten für die involvierten Subjekte benötigt.
Stattdessen verwendet \gls{fsgan} ein rekurrentes Netzwerk, das sowohl die Pose als auch die Gesichtsausdrücke kontinuierlich anpasst.
Ein wesentliches Element dieser Methode ist die Verwendung der Delaunay-Triangulation zur Interpolation zwischen verschiedenen Ansichten des Gesichts.
Diese Triangulation ermöglicht es, die Gesichtsausdrücke nahtlos und realistisch zu übertragen, selbst bei großen Pose- und Ausdrucksunterschieden.

Um verdeckte Gesichtsregionen zu rekonstruieren, setzt \gls{fsgan} auf ein Gesichtsvollständigungsnetzwerk, das fehlende Teile des Gesichts basierend auf den vorhandenen Segmentierungsdaten ergänzt.
Schließlich sorgt ein Gesichtsmischungsnetzwerk dafür, dass das rekonstruierten Gesicht harmonisch in das Zielbild eingefügt wird.
Diese Methode verwendet einen neuen Poisson-Blending-Verlust, der Poisson-Optimierung mit einem Wahrnehmungsverlust kombiniert, um eine nahtlose Integration zu gewährleisten.

Frühere Ansätze zum Reenactment stützten sich oft auf 3D-Modelle, die Gesichtslandmarken zur Steuerung der Gesichtsausdrücke nutzten.
Solche Modelle erforderten eine präzise Anpassung und waren anfällig für Fehler bei teilweise verdeckten Gesichtern.
Moderne Ansätze, wie \glspl{gan}, haben diese Herausforderungen durch die Nutzung tief lernender Netzwerke weitgehend überwunden, indem sie realistische und kohärente Ergebnisse liefern, ohne auf 3D-Modelle angewiesen zu sein.
\gls{fsgan} setzt neue Maßstäbe, indem es eine vollständig trainierbare, subjektsagnostische Lösung für Face Swapping und Reenactment bietet, die qualitativ und quantitativ überlegen ist\cite{face-swapping-and-reenactment}. % Kapitel

\end{document}
